\chapter{公式输入与手写样式}\label{ch:math}
\section{输入规则}
0.1.24 起可在公式中直接使用汉字，包括下标与上标，例如 \code{U_{模块}}、\code{U_{万}}、\code{温度}；分式、根式和矩阵同样支持。汉字使用已载入的真实字体，缺字继续定位到原文，不自动改用屏幕字体。\code{\text{模块}} 写法仍然支持。
Markdown 行内公式使用单美元符号；独立公式的起止双美元符号分别独占一行。代码区内的美元符号保持字面内容。
\par\noindent\begin{minipage}{\linewidth}
\begin{Source}
12.50 的一半为 $\frac{12.50}{2}$。

$$
\sqrt{12.50}
$$
\end{Source}
\end{minipage}\par\medskip
示例中的中文需要汉字资源。公式命令名不会逐字书写，但其显示结果需要实际字形，例如 \code{\alpha} 需要相应希腊字母轨迹，\code{\sin} 需要函数显示所用字母。内置数字不能代替完整数学字库。

\section{常用结构速查}
\par\noindent\begin{tabularx}{\linewidth}{@{}P{33mm}Y@{}}
\toprule
结构 & 输入示例\\\midrule
上下标 & \code{x_1^2}、\code{x^{12}}\\
分式 & \code{\frac{1}{2}}、\code{\dfrac{1}{2}}、\code{\tfrac{1}{2}}\\
根式 & \code{\sqrt{12}}、\code{\sqrt[3]{8}}\\
伸缩括号 & \code{\left(\frac{1}{2}\right)}\\
上下限 & \code{\sum_{i=1}^{n}}、\code{\int\limits_0^1}\\
文字与单位 & \code{\text{电压}}、\code{\mathrm{m^2}}\\
修饰符 & \code{\vec{x}}、\code{\overline{x}}、\code{\hat{x}}、\code{\dot{x}}\\
\bottomrule
\end{tabularx}

上下标超过一个字符时加花括号。\code{\left} 与 \code{\right} 要配对；不可见一侧可写点号。管道表格中的公式尽量用 \code{\vert} 或 \code{\Vert} 表示竖线，避免未转义的管道符被当作分列标记。

公式作为整体参与换行与分页，不会自动在中间折行。超宽时，可改为显式多行结构、减小公式字号或扩大有效区域。上下标层级过深、字形过小也会报错，不会无限缩小到不可书写。

\section{在设置窗口调节结构}
\entry{设置 → 元素 → 公式设置…}
设置窗口按结构分类：字号与字形、间距与数学轴、上下标、分式、根式、积分与大运算符、矩阵与向量、括号与定界符、多行对齐、公式编号、修饰与向量。

左侧选类别，在上方紧凑结构图旁的小输入框直接调节参数。局部短箭头标明参数对应的尺寸，聚焦时突出显示；窗口放大不会拉开字符与控件。示意图使用正体字形，保持固定比例，数值变化请以下方预览为准。下方实际字形预览使用当前项目笔迹，有效数值输入后约 0.1 秒刷新，无需按回车或离开输入框。预览倍率固定，修改字号不会被自动缩放抵消；可在倍率下拉框主动缩放，较大公式可滚动查看。超范围输入显示原因，不继续显示旧预览。

示例输入只填公式内容，不加美元符号；可切换行内与独立预览。示例下拉框提供小矩阵、带框线数组、连分式等试排内容，切换分类会保留各类自定义示例。示例不会自动插入正文；缺字仍显示红框。

默认公式字号为 6 mm，与表格数字高度独立。字号可在 mm 与 pt 之间联动显示；百分比参数以当前层公式字号为参照，例如分子、分母和上下标可能已经缩小。将鼠标悬停在参数旁的\ui{i}上，可查看参数含义和按主字号估算的毫米尺寸；点击图标也可打开提示。分类和预览标题旁的\ui{i}提供结构规则与资源说明。估算不能代替实际嵌套尺寸。勾选\ui{显示辅助线}可查看灰色基线和蓝色数学轴；所有尺寸箭头和辅助线均不进入机器文件。

参数框旁的回转箭头只恢复单项；\ui{当前分类恢复默认}重置该类尺寸、来源及该类对齐选项。\ui{应用设置}应用到文档并随项目保存，\ui{取消}保留原设置。本版没有逐条公式的独立样式覆盖，调整时应同时查看文档中其他公式。

普通上下标不再提供“净间距”手动项。程序按实际笔迹边界计入笔宽，自动保持水平留白及上下标之间的必要间隔；上移、下移设置仅控制最小位移，必要时会自动避让。分类标题旁的\ui{i}说明这一行为。

在\ui{积分与大运算符}中，\ui{上下限字号}独立于普通上下标字号，作用于求和、连乘、积分等的右侧或堆叠上下限。比例以当前层公式字号为基准，默认 70\%，可调范围为 50\%--100\%；上下限内部的嵌套上下标仍沿用普通上下标比例。

旧项目缺少独立上下限比例时，继承原有上下标比例后独立保存。普通上下标按新的自动留白规则重新排版，旧净间距值不再控制几何；打开旧文档后请检查分页与公式位置。

\section{使用录制的减号、根号与括号}
三类构件可以分别选择\ui{几何构件}或\ui{录制字形（缩放适配）}。自 0.1.27 起，新文档分数线默认使用字体中的数学减号横向拉伸，保留纵向笔迹起伏；根号和括号仍默认几何构件。旧文档保持已保存来源。缺少减号时可补录，或明确选择几何构件，程序不会静默替换。
\begin{enumerate}
\item 在\ui{分式}中选择录制来源，用数学减号的轨迹按内容宽度延长分数线。
\item 在\ui{根式}中选择录制根号，程序适配钩部、竖部与上横线；根内内容变长时主要延长横线。
\item 在\ui{括号与定界符}中选择录制来源，对相应括号按内容高度与构件宽度缩放。矩阵、分段和编号括号也受这一选择影响。
\end{enumerate}
\shot{settings-fraction.png}{0.94}{紧凑分式设置：短箭头连接附近输入框，说明收纳于 i 图标，下方核对实际笔迹。}

数学减号是 U+2212，不等同于正文中的 ASCII 连字符。缺样、退化轨迹或不能可靠识别的根号会报错。可补充合适样本，或明确切回几何模式；程序不会悄悄换来源。适配只作用于派生结构，不改原始采集数据。

\section{矩阵、分段与多行推导}
普通矩阵支持 \code{matrix}、\code{pmatrix}、\code{bmatrix} 等环境；\code{cases} 用于两列分段。矩阵以 \code{&} 分列，以双反斜线换行；普通矩阵各行列数应一致，最多 8 行、8 列。
\par\noindent\begin{minipage}{\linewidth}
\begin{Source}
$$
\begin{pmatrix}
1&2\\
3&4
\end{pmatrix}
$$
\end{Source}
\end{minipage}\par\medskip
多行推导可用 \code{aligned}，对齐点放在等号之前。下例需要加号、等号和数学减号等相应资源。
\par\noindent\begin{minipage}{\linewidth}
\begin{Source}
$$
\begin{aligned}
12+3&=15\\
15-2&=13
\end{aligned}
$$
\end{Source}
\end{minipage}\par\medskip
\code{gathered} 逐行居中；\code{split} 用于一个展示公式中的共同对齐点。多行结构仍整体分页，不跨页拆行。需要更多可独立分页的步骤时，应拆成多个独立公式。

\section{编号与引用}
普通独立公式默认不编号。\code{equation} 为整块生成编号，\code{align}、\code{gather} 为每行生成编号；对应星号环境不自动编号。
\par\noindent\begin{minipage}{\linewidth}
\begin{Source}
$$
\begin{equation}
\frac{1}{2}\label{eq:half}
\end{equation}
$$

$\eqref{eq:half}$
\end{Source}
\end{minipage}\par\medskip
\code{\tag{12.3}} 指定可见标签，\code{\tag*{12.3}} 不带括号；\code{\notag} 取消该行编号。\code{\ref} 显示标签文字，\code{\eqref} 显示带括号的引用。引用命令要放在数学定界符内。

\code{\label} 的键全篇唯一，引用可以出现在被引公式之前。重名、缺失目标或在无编号行登记标签会报错。插入前面的公式后，自动编号与引用会更新；不需要手动修改正文数字。

\section{与完整 LaTeX 的区别}
编辑器接受受控数学语法，不编译完整 \code{.tex} 文档。不要粘贴 \code{\documentclass}、\code{\usepackage}、宏定义或文件输入命令。\code{array} 支持 \code{l/c/r} 列及受控框线，但不支持任意列格式或 \code{\multicolumn}。未知命令会定位报错，保留原文供修改。

数学校准提醒表示当前参考线还需实际样张复核；它不同于“缺少字形”的阻断错误。资源齐全后，也应检查上下标、分数线和括号在实际笔宽下是否清楚。
